\newpage\pagestyle{conclusions}

\unnumberedchapter{Conclusioni e sviluppi futuri}
\label{chap:conclusions}

\unnumberedsection{Sintesi del lavoro}

Il lavoro è partito da una domanda operativa: se sia possibile selezionare le regole rilevanti in
un insieme di metadati impiegato per l'imputazione di valori mancanti, riducendo il costo
computazionale senza degradare la qualità delle imputazioni. Il sistema studiato, \renuver{},
riceve le regole come input esterno, e la selezione è stata quindi affrontata nella forma di
\emph{preprocessing}: riscrivere il file delle regole, senza toccare l'eseguibile.

La risposta a quella domanda si è articolata in tre passaggi. Il primo è stato la ricostruzione
del comportamento del sistema, per il quale non erano disponibili i sorgenti: l'analisi del
bytecode ha mostrato che una regola partecipa a due procedure con requisiti opposti --- la
generazione dei candidati e la verifica dei veti --- e che nulla, nel formato del file, distingue
i due ruoli. Il secondo passaggio è stato la derivazione delle condizioni esatte di potatura che
ne discendono, e del risultato negativo che ne consegue. Il terzo è stato la misurazione
sistematica dell'effetto delle strategie sviluppate, condotta su tre famiglie di dataset per un
totale di \num{1618} esecuzioni in configurazione potata.

Alla domanda iniziale il lavoro risponde su tre piani distinti, che conviene richiamare in chiusura
senza ripetere l'analisi: la \emph{riduzione garantita} non può superare il \SI{20.9}{\percent}
sui dataset esaminati, e il limite è del criterio e non dell'implementazione
(\autoref{chap:theory}); la \emph{qualità} può migliorare, ma soltanto dove la composizione delle
imputazioni lascia spazio al miglioramento e dove la copertura non cala, come mostra il confronto
fra Bridges e le altre famiglie (\autoref{chap:results}); il \emph{tempo} non segue la riduzione,
perché dipende da quali regole sopravvivono e non da quante (\autoref{sec:res-cost}). Le tre
risposte sono riassunte nella tabella della \autoref{sec:research-question} e discusse nel
\autoref{chap:discussion}, dove sono raccolti anche gli esiti delle ipotesi e i limiti
dell'evidenza.

\unnumberedsection{Possibili sviluppi}\label{sec:future-work}

\paragraph{Riscrivere le regole invece di rimuoverle.}
Il teorema negativo delimita le trasformazioni che \emph{eliminano} regole operando sul solo file.
Una trasformazione che riscriva le regole --- restringendo o allargando i lati sinistri, correggendo
le soglie, fondendo più regole in una --- non ricade in quel limite, e la condizione \Vtwo{} mostra
del resto che una regola può essere eliminata quando il suo contenuto informativo è conservato da
un'altra: la stessa conservazione potrebbe essere cercata per via sintattica anziché per
soppressione. È la direzione che più probabilmente consente di superare il tetto del
\SI{20.9}{\percent} senza rinunciare alla garanzia.

\paragraph{Completezza della condizione di copertura.}
La condizione \Vtwo{} è sufficiente ma non necessaria: stabilire se esistano altre classi di
regole il cui veto è implicato da una regola sopravvissuta --- e se siano decidibili con
l'informazione disponibile --- amplierebbe la portata quantitativa del risultato senza cambiarne la
natura. Il problema è aperto e riguarda la teoria delle dipendenze rilassate, non
l'implementazione.

\paragraph{Verificare il meccanismo del costo.}
L'ipotesi che il rallentamento osservato dipenda dalla perdita della scorciatoia esatta non è
verificata. Isolare il percorso della scorciatoia con la strumentazione e contare quante celle lo
impiegano, prima e dopo la potatura, è un esperimento circoscritto che trasformerebbe una
spiegazione plausibile in un dato, e che consentirebbe di prevedere il tempo di una strategia
oltre che la sua qualità.

\paragraph{Ottimizzare il percorso dei veti.}
La caratterizzazione del costo indica una direzione precisa e diversa da quella seguita finora.
La verifica dei veti esegue \num{354809} scansioni del dataset su \num{71} tuple distinte: la
circostanza suggerisce che il verdetto per coppia (regola, tupla) sia memoizzabile. Un
\emph{indice} sulle colonne del dataset, o un ordinamento delle regole che incontri presto quelle
che violano, sono direzioni alternative. Si tratta di interventi che richiedono la modifica del
sistema, ora disponibile come possibilità ma non ancora sperimentata su questo fronte.

\paragraph{Un teorema lato generazione.}
La condizione \Vtwo{} copre il ruolo di veto. Per il ruolo generativo non si dispone di una
condizione sufficiente analoga: la strategia \str{score-argmin} è esatta rispetto alla generazione
ma calcola l'argomento del minimo sull'insieme \emph{originale}, e la sua esattezza decade quando
altre regole sono state rimosse. Una condizione che garantisca l'invarianza dell'argomento del
minimo \emph{dopo} altre rimozioni colmerebbe la lacuna e renderebbe componibili i due criteri.

\paragraph{Misurare l'esercizio effettivo dei veti.}
La politica di tutela dei veti impiega un indice di rischio che, come si è mostrato, è
anti-correlato con la rilevanza. La sua sostituzione con una misura del numero di veti
effettivamente esercitati a runtime --- ottenibile dalla strumentazione già disponibile ---
consentirebbe di sostituire la stima con un dato.

\paragraph{Validare il criterio diagnostico su una famiglia indipendente.}
Il criterio della \autoref{sec:diagnostic} è stato derivato su tre famiglie e non è stato
sottoposto a validazione prospettica. Una famiglia con quota esatta intermedia e poche regole
costituirebbe il test decisivo.

Il risultato di questo lavoro non è una strategia di potatura che funzioni ovunque: è la
determinazione di \emph{che cosa} si può garantire, di \emph{quanto} si può guadagnare
abbandonando la garanzia, e del \emph{meccanismo} che decide quale dei due regimi si applichi a un
dataset dato. In un ambito in cui la letteratura presenta prevalentemente risultati positivi su
configurazioni favorevoli, la determinazione di un limite dimostrabile e la sua quantificazione
hanno un valore che non dipende dal segno del risultato.
